\begin{frame}{Design focus: LOUIS-guided}
\begin{columns}[T,onlytextwidth]
\begin{column}{.35\textwidth}
\begin{alertblock}{\faKey~Design focus}
\begin{itemize}
    \item limited set of competences (LOUIS)
    \item sustained industry context
    \item iterative development toward capstone
\end{itemize}
\end{alertblock}
\end{column}

\begin{column}{.61\textwidth}

\begin{block}{LOUIS framework}
\small
\textbf{Learning Outcomes in University for Impact on Society}

Focus on a small set of competences (max 3)

\end{block}
\vspace{0.3em}

\centering
\resizebox{\linewidth}{!}{\tikzset{
  compbox/.style={
    draw, rounded corners=5pt, align=center,
    text width=30mm, minimum height=8mm,
    font=\scriptsize, inner sep=1pt
  },
  highlight/.style={
    compbox,
    fill=green!20, draw=green!60!black, very thick
  }
}

\begin{tikzpicture}
  \matrix[matrix of nodes,
          nodes={compbox},
          column sep=2mm, row sep=2mm] {
    {Civic engagement} &
    {Creative thinking} &
    {Critical thinking} &
    {Ethical reasoning} \\
    {Global learning} &
    |[highlight]| {Information literacy} &
    {Inquiry and analysis} &
    {Integrative learning} \\
    {Intercultural knowledge \& competence} &
    {Foundations for life-long learning} &
    {Oral communication} &
    |[highlight]| {Problem solving} \\
    {Quantitative literacy} &
    {Reading} &
    |[highlight]| {Teamwork} &
    {Written communication} \\
  };
\end{tikzpicture}}

\end{column}
\end{columns}
\end{frame}
\note{
For this redesign, I attended a workshop on the LOUIS framework --  
\emph{Learning Outcomes in University for Impact on Society}.  

LOUIS is part of the Aurora University Network,
and is intended to strengthen broader academic and personal competences  
alongside subject knowledge.  

A key takeaway from that workshop was to avoid spreading too thin --  
instead, focus on a \emph{small number of competences}  
and make expectations explicit.  

So this redesign is directly informed by that work.  

LOUIS helped translate high-level competence goals  
into concrete learning outcomes  
that students can understand and act on.  

}
\begin{frame}{Pedagogical structure}
\centering\resizebox{.8\linewidth}{!}{\begin{tikzpicture}
    [scale=0.62, transform shape,
    node distance=65mm,
    inner sep=3pt,
    auto,
    block/.style={
        draw,
        thick,
        fill=green!40,
        text width=30mm,
        align=center,
        minimum height=10mm
    },
    line/.style={
        draw, thick, ->, shorten >=2pt
    },
    dashedline/.style={
        draw, thick, ->, dashed, shorten >=2pt, color=black!50
    },
    week/.style={anchor=south west, text width=45mm},
    red/.style={block, draw=red, fill=red!40},
    orange/.style={block, draw=orange, fill=orange!40},
    yellow/.style={block, draw=yellow, fill=yellow!40},
    green/.style={block, draw=green, fill=green!40},
    violet/.style={block, draw=violet, fill=violet!40},
    blue/.style={block, draw=blue!40, fill=blue!20},
    agile/.style={block, draw=gray, fill=gray!10,inner sep=3mm}
    ]

    % Samantektarverkefni
    \node[week] (themecapstone) at (0,-75mm) {Week 11--14\\ Capstone Project};
    \node[blue, right of=themecapstone, xshift=70mm, label=below:\textbf{Project}] (capstone) {Team (tCAP)};
    % Endurgjöf milli teyma
    \draw [line] (capstone) edge[loop above] (capstone);
    \node[blue, right of=capstone, label=below:\textbf{Reflection}] (capref) {Student (iCAP)};
    \draw[line] (capstone) to (capref);

    % Litur/lotur
    \foreach \week/\theme/\color/\pos in {1--2/Data Ethics/red/1, 3--4 /Large Language Models/orange/2, 5--6/Supervised Learning/yellow/3, 7--8 /Clustering/green/4, 9--10/Process Mining/violet/5}{
        \node[week] (week\pos) at (0,-\pos*10mm) {Week \week\\\theme};
        \node[\color, right of=week\pos, xshift=-50+10*\pos,label=below:\textbf{Preparation}] (RAT\pos)  {Student (iRAT)\\Team (tRAT)};
        \node[\color, right of=RAT\pos,label=below:\textbf{Project}] (APP\pos) {Team (tAPP)};
        \node[\color, right of=APP\pos,label=below:\textbf{Reflection}] (REF\pos) {Student (iREF)};

        % Línur milli eininga
        \path [line] (RAT\pos) -- (APP\pos);
        \path [line] (APP\pos) -- (REF\pos);

        % Jafningjamat innan Teams
        \draw [line] (APP\pos) edge[loop above] (APP\pos);
        
        % Tengingar við samantektarverkefni
        \path [dashedline] (APP\pos.east) to[out=-20, in=70] (capstone.north east);
        \path [dashedline] (REF\pos.east) to[out=-30, in=40] (capstone.east);
    }
    \begin{scope}[on background layer]
    \node[agile,fit=(RAT1) (RAT5),label=above:{\textbf{Sprint planning}}] () {};
    \node[agile,fit=(APP1) (APP5),label=above:{\textbf{Sprint}}] () {};
    \node[agile,fit=(REF1) (REF5),label=above:{\textbf{Backlog}}] () {};
    \end{scope}
\end{tikzpicture}
}
\begin{columns}[T,onlytextwidth]
    \begin{column}{.45\textwidth}
        \begin{itemize}
            \item[\faUserPlus] Team-Based Learning (TBL):
            \begin{itemize}
                \item[\faUser] iRAT (individual readiness)
                \item[\faUsers] tRAT (team readiness)
                \item[\faUsers] tAPP (team application)
                \item[\faUser] iREF (individual reflection)
            \end{itemize}
        \end{itemize}        
    \end{column}
    \begin{column}{.55\textwidth}
        \begin{itemize}            
            \item[\faShuffle] Flipped classroom
            \item[\faRepeat] Iterative modules leading into capstone
        \end{itemize}                
    \end{column}
\end{columns}
\end{frame}
\note{
The course combines \emph{TBL} with \emph{iterative} project work via \emph{PBL}  

Students prepare individually before class,  
and class time is used for discussion and guided project work.  (flipped classroom)

If we translate this into an Agile perspective:  
\begin{itemize}
    \item The preparation and initial team phase  
function like \emph{sprint planning}.  
    \item The subsequent sessions correspond to the \emph{sprint itself},  
where teams develop and refine their work.  
    \item And the individual reflection acts as a \emph{retrospective},  
capturing what should be improved next.  
    \item The capstone then brings this together.  
Rather than starting a new cycle,  
students revisit and refine previous work --  
which is closer to what in Agile would be called  
\emph{iteration on an evolving product}, not isolated sprints.  

    \item The final course reflection then acts as a broader retrospective.

\end{itemize}
}
